\documentclass{article}
\usepackage[utf8]{inputenc}
\usepackage{amsmath,amssymb}
\usepackage[most]{tcolorbox}
\usepackage{geometry}
\geometry{margin=2.5cm}

\newcommand{\step}[1]{\par\noindent\hangindent=3.5em\hangafter=1%
  \makebox[3.5em][l]{\textbf{#1.}}}
\newcommand{\steptag}[1]{\hfill{\small\textsf{[#1]}}}

\newtcolorbox{proofbox}[1]{
  colback=gray!5, colframe=gray!60,
  fonttitle=\bfseries, title={#1},
  breakable, enhanced,
  left=4pt, right=4pt, top=4pt, bottom=4pt,
  before skip=10pt, after skip=10pt
}

\begin{document}

\begin{proofbox}{For unital positive Phi on B(H)\_sa with ||Phi\textasciicircum{}2-Phi|| <= ...}
\step{1} For unital positive Phi on B(H)\_sa with ||Phi\textasciicircum{}2-Phi|| <= eta <= eta\_0, the near-fixed algebra (A=Im P, •, 1, A∩B(H)\_+) is an eps-JB algebra with eps <= C sqrt(eta), with universal dimension-free constants and no use of complete positivity. \steptag{validated} \\[6pt]
\step{1.1} ORDER-UNIT STRUCTURE (exact). By GT-orderunit, the near-fixed algebra (A=Im P, 1, A\_+=A cap B(H)\_+) is a finite-dimensional real order-unit space, and its order-unit norm ||a||\_ou=inf\{t>0: -t1<=a<=t1\} coincides EXACTLY with the operator norm on A. By near\_fixed\_algebra (def-near-fixed-algebra) A=Im P is well-defined: P=theta(2Phi-1) is the spectral idempotent (GT-Pprops gives P\textasciicircum{}2=P, P(1)=1, so 1 in A and A is a real subspace of B(H)\_sa). Thus the def-eps-jb-algebra ambient requirement -- a finite-dim real order-unit space with order-unit norm -- holds EXACTLY (zero eta-defect in order/unit/norm). \steptag{validated} \\[4pt]
\step{1.2} UNIT LAW AND COMMUTATIVITY (exact). By GT-easy, the projected product a•b=P(a o b) on A satisfies the unit law 1•a=a and commutativity a•b=b•a EXACTLY (no eta error). These are precisely the two exact identities required by def-eps-jb-algebra (commutativity a o b=b o a and unit law 1 o a=a hold exactly), with the product symbol o of the definition realised by • on A. \steptag{validated} \\[4pt]
\step{1.3} AXIOM JB1 (approximate submultiplicativity). By GT-easy, ||a•b|| <= (1+C\_1 eta)||a|| ||b|| for all a,b in A, with C\_1 a universal dimension-free constant. This delivers def-eps-jb-algebra (JB1) ||a o b|| <= (1+eps)||a|| ||b|| with the O(eta) defect C\_1 eta (to be absorbed into eps=C sqrt(eta) at node 1.7). \steptag{validated} \\[4pt]
\step{1.4} AXIOM JB2 (lower square-norm bound). By GT-easy, ||a•a|| >= (1-C\_2 eta)||a||\textasciicircum{}2 for all a in A, with C\_2 universal dimension-free. This delivers def-eps-jb-algebra (JB2) ||a o a|| >= (1-eps)||a||\textasciicircum{}2 with the O(eta) defect C\_2 eta (absorbed at node 1.7). \steptag{validated} \\[4pt]
\step{1.5} AXIOM JB3 (approximate positivity of squares). By GT-easy, a•a >= -C\_3 eta ||a||\textasciicircum{}2 1 for all a in A, with C\_3 universal dimension-free. This delivers def-eps-jb-algebra (JB3) a o a >= -eps ||a||\textasciicircum{}2 1 with the O(eta) defect C\_3 eta (absorbed at node 1.7). \steptag{validated} \\[4pt]
\step{1.6} AXIOM JB4 (approximate Jordan identity). By GT-jordan, ||((a•a)•b)•a - (a•a)•(b•a)|| <= C\_4 sqrt(eta) ||a||\textasciicircum{}3 ||b|| for all a,b in A, with C\_4 universal dimension-free. This delivers def-eps-jb-algebra (JB4) ||((a o a) o b) o a - (a o a) o (b o a)|| <= eps ||a||\textasciicircum{}3 ||b|| with the O(sqrt(eta)) defect C\_4 sqrt(eta) (this is the dominant term setting eps; see node 1.7). \steptag{validated} \\[4pt]
\step{1.7} ABSORPTION: fix eta\_0 < 1/4 (e.g. eta\_0=1/8), strictly below 1/4 — the BINDING constraint from GT-Pprops (lem-P-properties requires eta<=eta\_0<1/4 for the binomial series R=(S\textasciicircum{}2)\textasciicircum{}\{-1/2\}=(1-4(Phi-Phi\textasciicircum{}2))\textasciicircum{}\{-1/2\} to converge, i.e. 4 eta\_0<1). Since eta\_0<1/4<1, for 0<=eta<=eta\_0 one has eta<=1 => eta\textasciicircum{}2<=eta => eta<=sqrt(eta). Hence the O(eta) defects (JB1-JB3: C\_1,C\_2,C\_3 eta) and the O(sqrt eta) defect (JB4: C\_4 sqrt eta) combine into eps=C sqrt(eta), C=max(C\_1,C\_2,C\_3,C\_4) a FINITE universal constant; final exponent sqrt(eta). \steptag{validated} \\[4pt]
\step{1.8} Constants universal \& dimension-free: with eta\_0 < 1/4 FIXED (strict), the GT-Pprops constant C(eta\_0)=sum\_\{n>=1\}|a\_n|(4 eta\_0)\textasciicircum{}\{n-1\} (a\_n the binomial coefficients of x\textasciicircum{}\{-1/2\}) is FINITE because 4 eta\_0<1 (inside the radius of convergence rho=1); hence every C\_i built from it is finite, universal, and dimension-free (no dependence on dim H or rank). At the EXCLUDED boundary eta\_0=1/4 (4 eta\_0=1) the series diverges — which is precisely why eta\_0<1/4 must be strict. AND no complete positivity is used anywhere: the binomial P, Jordan-Schwarz, and the one-hole estimates each use only positivity+unitality (CP-free per their validated workspaces); the assembly never invokes an amplification id\_n (x) Phi. \steptag{validated} \\[4pt]
\end{proofbox}

\end{document}
